\chapter{Итог части II: точная граница дефектно-транспортной программы}

Вторая часть Тома V началась с отказа от статического выбора готового
состояния. Вместо него первичным был принят локальный переход, а частица
рассматривалась как устойчивый дефект композиции таких переходов. Эта
перестановка дала содержательную цепочку результатов, но одновременно
позволила точно увидеть границу текущей архитектуры.

Настоящая глава не вводит новый механизм. Она замораживает доказанное,
отделяет его от условных интерпретаций и формулирует единственное условие,
при котором дефектная ветвь может быть научно переоткрыта.

\section{Что удалось построить}

Положительное ядро части II образует следующую последовательность:
\begin{equation}
 \boxed{
 \begin{gathered}
 \text{типизированный переход}
 \longrightarrow
 \text{локальный перенос}
 \\
 \longrightarrow
 \text{голономный характер}
 \longrightarrow
 \text{самопорождающийся дефект}
 \\
 \longrightarrow
 \text{точечный проекторный заряд}
 \longrightarrow
 \text{условный спинорный индекс}.
 \end{gathered}}
 \label{eq:v5-part-conclusion-chain}
\end{equation}
Каждая сплошная стрелка в этой схеме имеет конкретное содержание.

\begin{enumerate}
 \item Носитель $M_{20\times15}(\mathbb C)$ допускает чтение не только как
 пространство состояний, но и как пространство элементарных переходов.
 \item Локальный унитарный шаг имеет дираковский непрерывный предел и общую
 предельную скорость.
 \item Цикл порядка четыре даёт точную задержку
 $1,i,-1,-i,1$, но не вечное хранение энергии.
 \item Нелинейный функционал с двумя вакуумами сам порождает кинк и
 локализованную дираковскую нулевую моду.
 \item В трёх пространственных измерениях кинк является стенкой; правильную
 точечную коразмерность даёт семейный проектор
 $P=\mathbf n\mathbf n^T$ с вакуумным пространством $\mathbb{RP}^2$.
 \item Проекторный ёж имеет заряд один, а спинорный подъём
 $\mathbf n\cdot\boldsymbol\sigma$ имеет собственную линию с $c_1=1$.
 \item Голономия, семейный проектор, лептонный блок и Хиггс дают условную
 лестницу внутреннего чтения $300\to45\to15\to2\to1$.
\end{enumerate}

Это уже больше, чем философская метафора. Построены язык, кинематика,
несколько точных топологических инвариантов и явные критерии провала.
Однако эта цепочка ещё не является моделью элементарной частицы.

\section{Где именно обрывается вывод}

Текущая архитектура не выводит два обязательных объекта.

\subsection{Пространственная семейная связность}

Для проекторного ежа чистый градиентный член даёт
\begin{equation}
 \Tr(\partial_iP\,\partial_iP)=\frac4{r^2},
 \qquad dE=8\pi\,dr.
\end{equation}
Следовательно, глобальная энергия линейно расходится. Конечная энергия
требует пространственной связности $A_i$ и ковариантной производной
\begin{equation}
 D_iP=\partial_iP+[A_i,P],
 \label{eq:v5-part-conclusion-covariant-projector}
\end{equation}
либо другого заранее выведенного стабилизатора. Имеющаяся операторная
семейная связность пока не доказана как физическое поле на пространстве.

\subsection{Ориентированный спинорный подъём}

Проектор $P$ хранит ось, но забывает её ориентацию:
\begin{equation}
 P(\mathbf n)=P(-\mathbf n).
\end{equation}
Прямая масса $3P-I_3$ имеет нулевой первый класс Черна. Единичный индекс
возникает только для $\mathbf n\cdot\boldsymbol\sigma$, меняющего знак при
$\mathbf n\mapsto-\mathbf n$. Поэтому нужна корневая или спинорная линия,
функториально восстанавливающая потерянный знак. Фаза порядка четыре
$h=\pm i$ подходит по алгебраическому типу, но её обязательное назначение
пакету $H_{15}$ не выведено.

\section{Почему это одна задача, а не две подгонки}

Недопустимо независимо добавить к модели калибровочное поле для энергии и
корневую линию для индекса. Тогда две трудности были бы устранены двумя
ручными расширениями. Допустимое переоткрытие должно получить оба объекта
из одной родительской структуры.

Минимальный кандидат должен иметь схематический вид
\begin{equation}
 \begin{split}
 S_{\rm parent}[P,A,\mathcal L,\psi,H]
 =\int d^4x\,\bigg[
 &\frac1{4g^2}\Tr(F_{\mu\nu}F^{\mu\nu})
 +\frac{f^2}{2}\Tr(D_\mu P D^\mu P)
 +V(P)
 \\
 &+\overline\psi\bigl(i\gamma^\mu D_\mu-mQ_{1/2}(P,\mathcal L)\bigr)\psi
 +\mathcal R_H(P,H)\bigg].
 \end{split}
 \label{eq:v5-part-conclusion-parent-target}
\end{equation}
Здесь $\mathcal L$ обозначает данные двойного подъёма, а
$\mathcal R_H$ --- последующее физическое чтение через Хиггс. Формула
\eqref{eq:v5-part-conclusion-parent-target} является не предложенным
действием, а перечнем типов членов, которые будущий вывод обязан получить.

Коэффициенты $g,f,m$, потенциал, представление связности и связь
$Q_{1/2}$ не разрешается выбирать после сравнения с наблюдаемыми. Они
должны следовать из одного следа, одной меры и одной симметрии.

\begin{nogocriterion}[Запрет раздельного ремонта]
Дефектная ветвь не считается переоткрытой, если конечная энергия получена
добавлением одного нового сектора, а единичный индекс --- независимым
добавлением второго. Пространственная связность и спинорный подъём должны
быть двумя проявлениями одной заранее объявленной родительской структуры.
\end{nogocriterion}

\section{Статус ранних образов}

Ранние слова о ``свете внутри скрученного бублика'', ``четырёхкратном
замедлении'' и ``петле времени'' теперь получают ограниченное прочтение.

\begin{itemize}
 \item Свет не отождествляется с внутренним содержимым частицы. Сохраняется
 лишь абстрактный безмассовый перенос с общей предельной скоростью.
 \item Число четыре сохраняется как порядок внутренней фазы и задержки, а
 не как универсальное уменьшение физической скорости света.
 \item Скручивание сохраняется как потеря ориентации в проекторе и
 необходимость двойного подъёма, а не как буквальная поверхность Мёбиуса.
 \item Внутренняя временная задержка не обеспечивает вечность. Устойчивость
 должна следовать из топологического заряда, конечной энергии и спектральной
 щели.
\end{itemize}

Таким образом, ранняя интуиция не подтверждена буквально, но часть её
структуры пережила строгий отбор: цикл порядка четыре, неориентированная
ось, двойной подъём, проходящий поток и дефект вместо материального
контейнера.

\section{Финальная теорема части II}

\begin{theorem}[Граница дефектно-транспортной программы]
Часть II Тома V построила математически согласованный путь от первичного
перехода к точечному проекторному дефекту и условному единичному
спинорному индексу. Она не построила конечную элементарную частицу:
энергия глобального ежа расходится, а спинорный подъём не следует из
одного проектора. Переоткрытие допустимо только при выводе из единого
нормированного родительского функционала как пространственной семейной
связности, так и ориентированного двойного подъёма.
\end{theorem}

\begin{protocol}
\begin{enumerate}
 \item язык первичных переходов: \textbf{построен};
 \item локальный унитарный перенос: \textbf{построен};
 \item цикл порядка четыре: \textbf{построен как задержка, не удержание};
 \item самопорождающийся дефект: \textbf{построен в модели стены};
 \item точечная проекторная топология: \textbf{построена};
 \item конечная энергия точечного дефекта: \textbf{не получена};
 \item единичный спинорный индекс: \textbf{условно получен после подъёма};
 \item сам спинорный подъём: \textbf{не выведен};
 \item физическое чтение одной нейтринной линии: \textbf{условно};
 \item единый нормированный функционал: \textbf{не построен};
 \item часть II как аудит выбранного класса: \textbf{завершена};
 \item единая физическая теория: \textbf{не построена}.
\end{enumerate}
\end{protocol}

Автоматическое продолжение цепочки гейтов после этой главы запрещено.
Новая часть открывается только при предъявлении одного конкретного
родительского принципа, который одновременно порождает
\eqref{eq:v5-part-conclusion-covariant-projector} и спинорный подъём.

Машинный реестр статусов сохранён в
\path{s2t_v5_defect_transport_part_conclusion_gate_results.json}.